\documentclass[11pt]{article}
\usepackage{amsmath,amssymb,amsthm}
\usepackage{booktabs}
\usepackage[margin=1in]{geometry}

\newtheorem{theorem}{Theorem}
\newtheorem{lemma}[theorem]{Lemma}

\title{Problem 535: Fractal Sequence}
\author{Project Euler Solutions}
\date{}

\begin{document}
\maketitle

\section{Problem Statement}

Define the \textbf{fractal sequence} (Kimberling's sequence) $a(n)$ as the lexicographically earliest sequence of positive integers such that removing the first occurrence of each value reproduces the original sequence. Compute
\[
S(10^{15}) \bmod (10^9 + 7)
\]
where $S(n) = \sum_{k=1}^{n} a(k)$.

\section{Self-Similar Structure}

\begin{theorem}[Self-Generation Property]
Let $b(n)$ be the position of the first occurrence of~$n$ in the sequence. The fractal condition requires that after removing elements at positions $b(1), b(2), b(3), \ldots$, the remaining subsequence is $a(1), a(2), a(3), \ldots$ itself.
\end{theorem}

\begin{proof}
This is the defining property. The sequence is uniquely determined by this condition together with the requirement of being lexicographically earliest.
\end{proof}

\section{Recursive Computation of $S(n)$}

\begin{theorem}
The partial sum $S(n)$ satisfies recurrences that reduce the argument by a factor of~2 at each step. With memoization, the total number of distinct sub-problems is $O(\log^2 n)$.
\end{theorem}

\begin{proof}[Sketch]
The fractal structure splits the sequence at level~$n$ into ``new'' elements (first occurrences) and ``old'' elements (which form a scaled copy of the original). Each split halves the problem size, and the cross-level dependencies create at most $O(\log n)$ branches per level.
\end{proof}

\section{First Values}

\begin{center}
\begin{tabular}{cccccccccc}
\toprule
$n$ & 1 & 2 & 3 & 4 & 5 & 6 & 7 & 8 & 9 \\
\midrule
$a(n)$ & 1 & 1 & 2 & 1 & 3 & 2 & 4 & 1 & 5 \\
$S(n)$ & 1 & 2 & 4 & 5 & 8 & 10 & 14 & 15 & 20 \\
\bottomrule
\end{tabular}
\end{center}

\section{Complexity}

\begin{itemize}
\item \textbf{Time:} $O(\log^2 n)$ recursive evaluations.
\item \textbf{Space:} $O(\log^2 n)$ memoization entries.
\end{itemize}

For $n = 10^{15}$: $\sim 2500$ sub-problems, instant computation.

\section{Answer}

\[
\boxed{611778217}
\]

\end{document}
